\documentclass[10pt,a4paper]{amsart}
\usepackage[margin=19mm]{geometry}
\usepackage{amsmath,amssymb,mathtools}
\usepackage[scr=rsfs,cal=euler]{mathalfa}
\usepackage{xfrac}
\usepackage[unicode,hidelinks,bookmarks=false]{hyperref}
\newcommand{\ie}{i.e.\ }
\newcommand{\cf}{cf.\ }
\newcommand{\resp}{respectively\ }
\newcommand{\Subsection}{\subsection}
\providecommand{\nobreakdash}{\nobreak\hbox{-}}
\setlength{\emergencystretch}{2em}
\multlinegap0pt
\usepackage{bm}
\usepackage{bbm}
\usepackage{dsfont}
\usepackage{xspace}
\usepackage{cancel}
\usepackage{mathtools}
\usepackage{amsthm}
\makeatletter
\@ifundefined{thm}{\newtheorem{thm}{Thm}}{}
\@ifundefined{lem}{\newtheorem{lem}[thm]{Lemma}}{}
\makeatother
\newcommand{\IM}{ \mathop{\mathrm{Im}}  }
\newcommand{\RE}{ \mathop{\mathrm{Re}} }
\newcommand{\abs}[1]{\left\vert #1 \right\vert}
\title[An example of a space $L^{p(\cdot)}$ on which the Cauchy-Ler]{An example of a space $L^{p(\cdot)}$ on which the Cauchy-Leray-Fantappi\`{e} operator for complex ellipsoid is not bounded}
\author{Aleksandr Rotkevich}
\date{Source: arXiv:2511.05997v1 (CC BY 4.0)}
\begin{document}
\maketitle

\noindent\emph{Source and licence.} An example of a space $L^{p(\cdot)}$ on which the Cauchy-Leray-Fantappi\`{e} operator for complex ellipsoid is not bounded. Version arXiv:2511.05997v1 (CC BY 4.0), distributed under the Creative Commons Attribution 4.0 International licence, as recorded in the arXiv licence metadata for this exact version and shown on the abstract page https://arxiv.org/abs/2511.05997v1. Full rights notice: licenses/arxiv-2511.05997-CC-BY-4.0.txt.\par\smallskip

\noindent\textit{Editorial scope.} Counted unit: the Lem whose source label is \texttt{lm:EstOfKernel} in the article (source0.tex lines 266--276, source label \texttt{lm:EstOfKernel}), delivered together with the article's own complete original proof of it (source0.tex lines 278--404, one \texttt{proof} environment of 127 lines). No mathematics is written, repaired or re-derived: statement and proof are the article's own text line for line. Required context retained uncounted: none. Every definition, display and label of those lines is retained unchanged. Omitted from this package and not redistributed: 1--265, 277--277, 405--522. Nothing inside a retained excerpt is omitted or altered: every retained body is delivered exactly as the article's lines are written, so no omission receipt of any kind accompanies this package. Citations: the retained excerpts carry no citation command at all, so no bibliography entry of the article is transcribed and no reference is redistributed. Reference adaptation: none was needed and none was made; every reference inside the delivered unit resolves inside it, so no unresolved reference or citation is produced. Printed slips kept literal: no printed word, symbol, hypothesis or number of the counted statement or of its proof was altered. Layout and class: the article's own class is \texttt{elsarticle} with packages including lineno, hyperref, amssymb, amsmath; the wrapper is the lane's pinned \texttt{amsart} editorial wrapper at 10pt on A4, which restates the theorem-like environments the retained text needs and adds the article's own macro definitions that the retained text needs (\texttt{\textbackslash IM}, \texttt{\textbackslash RE}, \texttt{\textbackslash abs}), with extra wrapper packages (bm, bbm, dsfont, xspace, cancel, mathtools). The delivered numbering is the unit's own. Assets: no retained span contains a figure environment, a graphics-inclusion command, a PDF-page inclusion, a table or a TikZ picture; no image file is copied into this package, the package declares no asset dependency and no image file is redistributed with it. Licence: version arXiv:2511.05997v1 of this article is distributed under the Creative Commons Attribution 4.0 International licence, as recorded in the arXiv licence metadata for this exact version and shown on the abstract page https://arxiv.org/abs/2511.05997v1. A full rights notice is delivered at licenses/arxiv-2511.05997-CC-BY-4.0.txt. Characters: every retained excerpt is pure ASCII, so the delivered engine, XeLaTeX with its default Latin Modern text font, reports no missing character.
\begin{lem}[Estimate of the kernel on boundary patches] \label{lm:EstOfKernel}
	There exist parameters $0<\gamma_1<\gamma_2$, $0<\gamma_3<\gamma_4$ such that for sufficiently small $\alpha>0$ 
	\begin{equation}\label{lm:EstOfKernel_eq1}
		-\RE\left(\frac{1}{w(\xi,z)^2}\right) \gtrsim \frac{1}{\alpha^2}
	\end{equation}
	uniformly for $\xi\in W_\alpha$ and $z\in V_\alpha$. Moreover, we can choose parameters $0<\gamma_1<\gamma_2$, $0<\gamma_3<\gamma_4$ such that
	\begin{equation}\label{lm:EstOfKernel_eq2}
		-\RE\left(\frac{1}{w(\xi,z)^2}\right) > 0 , 
	\end{equation}
	for $\xi\in W_{\alpha_1}$ and $z\in V_{\alpha_2}$ if $\alpha_1, \alpha_2>0$ are small enough.
\end{lem}

\begin{proof} 
	Indeed, let $\xi=(\xi_1,\xi_2) =(r_1e^{i\theta_1},r_2e^{i\theta_2}) \in W_\alpha$ and $z = (s_1e^{i\varphi_1},s_2e^{i\varphi_2}) \in V_\alpha$. Then
	$$w(\xi,z) = m_1r_1^{2m_1-1}\left(r_1-s_1e^{i(\varphi_1-\theta_1)}\right) + m_2r_2^{2m_2-1}\left(r_2-s_2e^{i(\varphi_2-\theta_2)}\right)  =  A_1 + A_2.$$
	
	Since
	$$\abs{ r_1-s_1e^{i(\varphi_1-\theta_1)} }\leq \abs{r_1}+\abs{s_1} \lesssim \alpha^{1/(2m_1)}
	$$
	then
	$$\abs{A_1}\lesssim r_1^{2m_1-1}\alpha^{1/(2m_1)} \lesssim \alpha.$$
	To estimate $A_2$ notice that
	$$\abs{r_2-s_2}=(1-r_1^{2m_1})^{1/(2m_2)} - (1-s_1^{2m_1})^{1/(2m_2)} \lesssim \abs{ r_1^{2m_1} -s_1^{2m_1} } \lesssim \alpha.$$
	Hence, since $\varphi_2-\theta_2\in (2\alpha,4\alpha)$, we conclude 
	$$A_2 \lesssim \abs{r_2-s_2}+s_2\abs{1-e^{i(\varphi_2-\theta_2)} }\lesssim \alpha.$$
	Combing these estimates we see that 
	$$\abs{w(\xi,z)} \lesssim \alpha.$$
	
	The estimates above are valid for arbitrary $\gamma_j>0,$ $j=1,\ldots4.$ However, we also need to control the argument of $w(\xi,z).$ This will be achieved by the appropriate choice of parameters $\gamma_j,$ $j=1,\ldots4.$ Let first $\gamma_4=\frac{\gamma_1}{2}$ and $\gamma_{3} =\frac{\gamma_{4}}{2}$ and recall that
	$$ \gamma_1\alpha \leq r_1^{2m_1} \leq \gamma_2\alpha; \quad  \gamma_3\alpha \leq r_1^{2m_1} \leq \gamma_4\alpha ;$$
	$$r_2^{2m_2} = 1-r_1^{2m_1};\quad s_2^{2m_2} =  1-s_1^{2m_1}.$$

	
	The choice of $\gamma_4=\frac{\gamma_1}{2}$ ensures that $2s_1^{m_1}\leq  r_1^{m_1}<1$ and choosing $\alpha>0$ small enough we may assume that $2^{-1/(2m_2)}<r_2<s_2< 1.$
	
	
	Consider first the real part of $w(\xi,z)$:
	\begin{multline} \label{lm:EstOfKernel_REw}
		\RE w(\xi,z) = m_1 r_1^{2m_1-1} \left(r_1-s_1\cos(\varphi_1-\theta_1)\right)+\\
		m_2 r_2^{2m_2-1} \left(r_2-s_2\cos(\varphi_2-\theta_2)\right) = B_1+B_2.
	\end{multline}
	The first term 
	$$B_1\leq m_1 r_1^{2m_1} \leq m_1 \gamma_2 \alpha.$$
	For the the second term notice that if $\alpha>0$ is small enough
	\begin{equation} \label{eq2}
		0<s_2-r_2= \frac{ s_2^{2m_2}-r_2^{2m_2} }{ \sum\limits_{k=0}^{2m_2} s_2^k r_2^{2m_2-k} } \leq 
		\frac{ r_1^{2m_1} - s_1^{2m_1} }{ 2m_2 r_2^{2m_2} } \leq \frac{ \gamma_2 -\gamma_4 }{ 2 m_2 }  \alpha + O(\alpha^2)
		\leq\frac{ \gamma_2 }{ m_2 } \alpha.
	\end{equation}
	
	Since $$1-\cos(\varphi_2-\theta_2) \leq 8\alpha^2$$ then 
	$$|B_2|\leq m_2 \abs{r_2-s_2} + m_2 s_2(1-\cos(\varphi_2-\theta_2))\leq {2\gamma_2}\alpha$$
	if $\alpha>0$ is small enough.
	Finally,
	$$\RE{w(\xi,z)} \leq (2+m_1) \gamma_2\alpha.$$
	The real part is, indeed, positive. 
	Analogously to \eqref{eq2} we see that
	\begin{equation} \label{eq3}
	s_2-r_2= \frac{ s_2^{2m_2}-r_2^{2m_2} }{ \sum\limits_{k=0}^{2m_2} s_2^k r_2^{2m_2-k} } \leq \frac{\gamma_2-\gamma_4}{2m_2} \alpha + O(\alpha^2)
	\end{equation}
	Hence,
	\begin{multline*}
		\RE{w(\xi,z)} = m_1 r_1^{2m_1} \left(1-\frac{s_1}{r_1}\cos(\varphi_1-\theta_1)\right)+\\
		m_2 r_2^{2m_2-1}  \left(r_2-s_2\cos(\varphi_2-\theta_2)\right) \geq
		m_1 r_1^{2m_1} \left(1-\frac{s_1}{2r_1}\right) -
		m_2  \left(s_2-r_2\right) \geq\\
		 m_1\left(1-\frac{1}{4}\right)\gamma_1 \alpha -   \frac{\gamma_{2}-\gamma_4}{2} \alpha +O(\alpha^2)
		 \leq \left(\frac{3}{4} m_1-\frac{\gamma_2}{2\gamma_1}\right)\gamma_1\alpha
		  = c(\gamma_2/\gamma_2)\gamma_1\alpha
	\end{multline*}
 	when $\alpha>0$ is small enough. Since $m_1\geq1$ then 
	we can choose the relation $\gamma_2/\gamma_1>0$ such that $c(\gamma_1,\gamma_2)>0$, which ensures that 
	$\RE w(\xi,z) >0$. %when $\alpha>0$ is small enough. 
	
	Consider now the imaginary part:
	\begin{multline}\label{eq:est:ImW}
		-\IM{w(\xi,z)} = m_1 r_1^{2m_1-1} s_1 \sin(\varphi_1-\theta_1) + 
		m_2 r_2^{2m_2-1} s_2 \sin(\varphi_2-\theta_2) > \\
		m_2 r_2^{2m_2-1} s_2 \sin(\varphi_2-\theta_2)\geq  \frac{m_2}{\pi}(\varphi_2-\theta_2) \geq\frac{2m_2}{\pi} \alpha.
	\end{multline}
	
	Hence,
	$$0< -\cot\arg w(\xi,z) =-\frac{\RE w(\xi,z)}{\IM w(\xi,z)} \leq \frac{2+m_1}{2m_2} \pi \gamma_2 < \frac{2+m_1}{2m_2} \pi \gamma_1  < \frac{1}{\sqrt{3} }
	$$
	if we choose $\gamma_1>0$ small enough. 
	
	Since $\RE{w(\xi,z)}>0$ and $\IM {w(\xi,z)}<0$ this guarantees $-\frac{\pi}{2}<\arg w(\xi,z)<-\frac{2\pi}{3}.$ Hence, $-\cos(2\arg w(\xi,z))>\frac{1}{2}$ and
	$$
	-\RE \frac{1}{w(\xi,z)^2} =-\frac{\cos(2\arg w(\xi,z))}{\abs{w(\xi,z)}^2} \gtrsim \frac{1}{\alpha^2}
	$$
	To prove that 
	$$-\RE\left(\frac{1}{w(\xi,z)^2}\right) \geq 0 , $$
	for $\xi\in W_{\alpha_1}$ and $z\in V_{\alpha_2}$ if $\alpha_1, \alpha_2>0$ are small enough we show that 
	$$\abs{\RE{w(\xi,z)}} \leq C(\alpha_1+\alpha_2) \leq \abs{\IM w(\xi,z)}.$$
	The estimate \eqref{eq:est:ImW} implies that
	$$ \abs{\IM w(\xi,z)} \geq \frac{m_2}{\pi}(\alpha_1+\alpha_2).$$
	
	%Suppose now $\alpha_1,\alpha_2<\alpha_0.$ 
	To estimate real part recall the formula \eqref{lm:EstOfKernel_REw}. Now
%	$$ 
%	\abs{r_1-s_1}\leq (\gamma_{2}\alpha_1)^{1/(2m_1)}+(\gamma_4\alpha_2)^{1/(2m_1)} <  2\gamma_{2}^{1/(2m_1)}(\alpha_1+\alpha_2)^{1/(2m_1)} ;
%	$$
	\begin{multline*}
	\abs{B_1}=r_1^{2m_1-1}\abs{ r_1-s_1\cos(\varphi_1-\theta_1) } \leq r_1^{2m_1} + r_1^{2m_1-1} s_1 \leq \\	\left( \gamma_2\alpha_1+(\gamma_2\alpha_1)^{ 1-1/(2m_1) }(\gamma_{4}\alpha_2)^{ 1/(2m_1) }\right)\leq 2\gamma_2(\alpha_1+\alpha_2)
	\end{multline*}
	To estimate the second term $B_2$ consider
	$$ 
	1-\cos(\varphi_2-\theta_2) = O((\alpha_1+\alpha_2)^2);
	$$
	$$
	r_2=\left(1-r_1^{2m_1}\right)^{1/(2m_2)} = 1-\frac{r_1^{2m_1}}{2m_2} + O(r_1^{4m_1}) = 1-\frac{r_1^{2m_1}}{2m_2} + O((\alpha_1+\alpha_2)^2);
	$$
	$$
	s_2 = 1-\frac{s_1^{2m_1}}{2m_2} + O((\alpha_1+\alpha_2)^2);
	$$
	Hence,
	\begin{multline*}
		\abs{s_2-r_2} =  \frac{ \abs{s_1^{2m_1}-r_1^{2m_1}} }{  2m_2 } + O(\alpha_1+\alpha_2)^2 \leq \\
		\frac{\gamma_2 + \gamma_4}{2m_2}(\alpha_1+\alpha_2)+O(\alpha_1+\alpha_2)^2\leq \frac{\gamma_2}{m_2} (\alpha_1+\alpha_2)
	\end{multline*}
	and
	$$\abs{B_2}\leq m_2\abs{s_2-r_2}+m_2(1-\cos(\varphi_2-\theta_2))\leq 
	\gamma_2 (\alpha_1+\alpha_2) +O(\alpha_1+\alpha_2)^2\leq 2 \gamma_2 (\alpha_1+\alpha_2)
	$$
	if $\alpha_0>0$ is small enough. 	
	
	Finally,
 	$$
		\abs{\RE w(\xi,z) }\leq \abs{B_1}+\abs{B_2} \leq 2\left(m_1+1\right)\gamma_{2}(\alpha_1+\alpha_2)
	$$
	Choosing $\gamma_{2}>0$ such that
	$$ 2\left(m_1+1\right)\gamma_{2} < \frac{m_2}{\pi}$$
	we have 
	$$\abs{\RE w(\xi,z)} < \abs{\IM w(\xi,z)}$$
	and
	$$-\RE \frac{1}{w(\xi,z)^2} = \frac{(\IM{w(\xi,z)})^2-(\RE{w(\xi,z)})^2}{\abs{w(\xi,z)}^4} >0$$
	for $\xi\in W_{\alpha_1}$ and $z\in V_{\alpha_1}$ 
	which finishes the proof of the Lemma.\qed	
\end{proof}

\end{document}
